\PassOptionsToPackage{table,xcdraw}{xcolor}
\documentclass[11pt, a4paper]{article}

\usepackage[T1]{fontenc}
\usepackage[utf8]{inputenc}
\usepackage{tgpagella}
\usepackage[spanish, es-noshorthands]{babel}
\usepackage{amsmath, amssymb}
\usepackage[most]{tcolorbox}
\usepackage[american]{circuitikz}
\usepackage[margin=2.5cm]{geometry}
\usepackage{fancyhdr}

\pagestyle{fancy}
\fancyhf{}
\setlength{\headheight}{16pt}
\lhead{\textbf{UCB Sede Tarija} | Ing. Mecatrónica}
\rhead{IMT-221 | Diagnóstico S07-2}
\renewcommand{\headrulewidth}{0.4pt}
\definecolor{ucbblue}{RGB}{0, 51, 102}

\begin{document}

\begin{tcolorbox}[colback=ucbblue!5, colframe=ucbblue, title=\textbf{Diagnóstico Breve: Transferencia del Método de Pequeña Señal}]
    Tiempo estimado: 8-10 minutos[cite: 15]. Trabaja individualmente. El objetivo es comprobar que sabes \textbf{cómo iniciar y estructurar} la derivación, no que memorizaste una fórmula final[cite: 15].
\end{tcolorbox}

\section*{Evaluación de Topología BJT}
Se te presenta la siguiente topología de amplificador polarizado en región activa. Se asume que el punto de operación $I_{CQ}$ ya es conocido.

\begin{center}
\begin{circuitikz}[scale=1, transform shape, thick]
    \draw (0,0) node[npn](Q){};
    \draw (Q.C) -- (0,1.5) node[above]{$+V_{CC}$};
    \draw (Q.B) to[short, -o] (-1.5,0) node[left]{$v_i$};
    \draw (Q.E) to[R, l_={$R_E$}] (0,-2.5) node[ground]{};
    \draw (Q.E) to[short, *-o] (1.5,-0.77) node[right]{$v_o$};
\end{circuitikz}
\end{center}

\textbf{Completa el esqueleto de derivación explícitamente:}
\begin{enumerate}
    \item \textbf{Identificación de terminales:} ¿Por qué terminal entra la señal, por cuál sale, y cuál será conectado a Tierra AC? \vspace{1cm}
    \item \textbf{Parámetros intrínsecos:} Escribe las fórmulas para calcular $g_m$ y $r_e$ a partir del punto $Q$. \vspace{1cm}
    \item \textbf{Equivalente AC:} Dibuja el circuito reemplazando fuentes DC por Tierra AC y el transistor por su Modelo T simplificado (o Híbrido-$\pi$). \vspace{3.5cm}
    \item \textbf{Leyes de Kirchhoff:} Escribe la primera ecuación (KVL o KCL) que plantearías en tu circuito equivalente para relacionar $v_i$ con $v_o$. \vspace{1.5cm}
    \item \textbf{Supuestos:} Escribe al menos un supuesto técnico bajo el cual esta derivación es válida. \vspace{1cm}
\end{enumerate}

\end{document}
